\title{Group Sequence Policy Optimization}

\begin{abstract}
We present Group Sequence Policy Optimization (GSPO), a reinforcement learning method designed for training large language models that offers robustness, scalability, and strong performance. In contrast to existing approaches that rely on token-level importance ratios, GSPO employs a sequence-based importance ratio, implementing sequence-centric clipping, reward assignment, and optimization. Our findings indicate that GSPO outperforms the GRPO algorithm in both training efficiency and effectiveness, significantly enhances the stability of Mixture-of-Experts (MoE) RL training, and supports streamlining RL system architectures. The advantages of GSPO have played a key role in the notable advancements observed in the latest Qwen3 models.
\end{abstract}

\section{Introduction}
\label{sec:intro}

Reinforcement learning (RL) has become a central approach for advancing language models \citep{o1,dpsk-r1,qwq32b,qwen3}. Leveraging large-scale RL enables language models to solve complex tasks, including high-level mathematics and programming challenges, by fostering extended and intricate reasoning abilities.

To maximize the benefits of RL through increased computational resources, it is critical to ensure that training remains stable and resilient. Nonetheless, contemporary RL techniques—such as GRPO \citep{grpo}—suffer from pronounced instability when applied to the training of very large language models, frequently culminating in drastic and irreversible model failures \citep{qwen3, minimax-m1}. Such instability poses a significant barrier to advancing the potential of language models via sustained RL optimization.

This study demonstrates that GRPO’s instability arises chiefly from the inappropriate use and subsequent breakdown of importance sampling weights within its algorithmic framework. This flaw induces substantial variance during training, which intensifies as the response length grows and is exacerbated by the clipping procedure, eventually leading to the deterioration of the model.

To overcome these principal challenges, we introduce \textbf{Group Sequence Policy Optimization (GSPO)}, a novel reinforcement learning algorithm specifically designed for large language model training. GSPO’s major contribution is its principled formulation of the importance ratio, grounded in sequence likelihood \citep{click}, thereby adhering to the core tenets of importance sampling. Moreover, GSPO calculates normalized rewards by leveraging the advantage values across multiple responses for a given query, guaranteeing consistency between rewards assigned at the sequence level and the optimization objectives.

Through our empirical studies, we show that GSPO markedly outperforms GRPO in terms of training stability, efficiency, and overall effectiveness. Notably, GSPO addresses the stability issues commonly encountered in reinforcement learning with large Mixture-of-Experts (MoE) models by design, making additional stabilization methods unnecessary and indicating promise for streamlining RL systems. These strengths of GSPO have played a pivotal role in driving the impressive performance gains observed in the most recent Qwen3 models. We anticipate that GSPO will serve as a robust and extensible algorithmic base for future large-scale reinforcement learning endeavors in the context of language model training.

\section{Preliminaries}

\paragraph{Notation}
Throughout this work, we represent an autoregressive language model with parameters $\theta$ as the policy $\pi_\theta$. The variable $x$ signifies a query, and $\mathcal{D}$ stands for the set containing all queries. For a response $y$ associated with a query $x$, the conditional probability under $\pi_\theta$ is expressed as $\pi_\theta (y | x) = \prod_{t=1}^{|y|} \pi_\theta (y_t | x, y_{<t})$, where $|y|$ refers to the total number of tokens present in $y$. Each pair consisting of a query and its response, denoted by $(x, y)$, can be evaluated using a verifier $r$, which produces a reward value $r(x, y) \in [0, 1]$.

\paragraph{Proximal Policy Optimization (PPO)} Using samples generated from the old policy $\pi_{\theta_\text{old}}$, PPO \citep{ppo} constrains the policy update within a proximal region of the old policy through the clipping mechanism. Specifically, PPO employs the following objective for policy optimization (we omit the KL regularization term hereinafter for brevity, as it is not the focus of this paper):

\begin{align}
\mathcal{J}_\text{PPO}(\theta) = \mathbb{E}_{ x \sim \mathcal{D},\, y \sim \pi_{\theta_\text{old}}( \cdot | x) }
\left[ \frac{1}{|y|} \sum_{t=1}^{|y|} 
\min \left( w_{t}(\theta) \widehat{A}_{t},  \, \mathrm{clip} \left( w_{t}(\theta), 1 - {\varepsilon}, 1 + {\varepsilon}\right) \widehat{A}_{t} \right)
\right],
\end{align}

Here, the importance ratio for token $y_t$ is given by $
w_{t}(\theta) = \frac{ \pi_{\theta} (y_{t} | x, y_{<t}) }{ \pi_{\theta_\text{old}} (y_{t} | x,y_{<t})} $, with the advantage $\widehat{A}_{t}$ for $y_t$ being provided by an auxiliary value model, and $\varepsilon$ denoting the threshold for clipping the importance ratios.

A significant difficulty encountered when implementing PPO is its strong dependence on the value model. In particular, the value model is typically comparable in scale to the policy model, which substantially increases both memory requirements and computational demands. Additionally, the performance of the algorithm is contingent upon the accuracy of its value estimations. Obtaining dependable value estimates is already complex, and adapting the value model to accommodate longer outputs and increasingly intricate tasks amplifies this complexity further.

\paragraph{Group Relative Policy Optimization (GRPO)} GRPO \citep{grpo} bypasses the need for the value model by computing the relative advantage of each response within a group of responses to the same query. Specifically, GRPO optimizes the following objective:

\begin{align}
\label{equ:grpo}
\mathcal{J}_\text{GRPO}(\theta) = \mathbb{E}_{ x \sim \mathcal{D},\, \{y_i\}_{i=1}^G \sim \pi_{\theta_\text{old}}( \cdot | x) }
\left[ \frac{1}{G} \sum_{i=1}^{G} \frac{1}{|y_i|} \sum_{t=1}^{|y_i|} 
\min \left( w_{i,t}(\theta) \widehat{A}_{i,t},  \, \mathrm{clip} \left( w_{i,t}(\theta), 1 - {\varepsilon}, 1 + {\varepsilon}\right) \widehat{A}_{i,t} \right)
\right],
\end{align}

where $G$ is the number of generated responses to each query $x$ (i.e., the group size), and the importance ratio $w_{i,t}(\theta)$ and advantage $\widehat{A}_{i,t}$ of token $y_{i,t}$ are:

\begin{align}
    w_{i,t}(\theta)=\frac{ \pi_{\theta} (y_{i,t} | x, y_{i,<t}) }{ \pi_{\theta_\text{old}} (y_{i,t} | x,y_{i,<t})},\quad\ 
    \widehat{A}_{i,t} = \widehat{A}_{i} = \frac{r(x, y_i) - \mathrm{mean} \left( \{ r(x, y_i) \}_{i=1}^G \right) }{ \mathrm{std} \left( \{ r(x, y_i) \}_{i=1}^G \right) },
\end{align}

respectively, where all the tokens in $y_i$ share the same advantage as $\widehat{A}_{i}$.

\section{Motivation}
\label{sec:motivation}

As models become larger, increasingly sparse (as seen in Mixture-of-Experts models), and generate longer responses, efficient hardware usage during RL requires executing rollouts with substantial batch sizes. For enhanced sample efficiency, it is common to subdivide the large rollout batch into several mini-batches for gradient computation. However, this practice naturally leads to an off-policy training regime, since the samples $y$ originate from a previous policy $\pi_{\theta_\text{old}}$, not the present policy $\pi_\theta$ being improved. This circumstance underlies why mechanisms such as clipping in PPO and GRPO are necessary: they serve to exclude samples that are excessively “off-policy” from influencing gradient calculations.

While mechanisms like clipping aim to manage this off-policy discrepancy, we identify a more fundamental issue in GRPO: \textit{its objective is ill-posed}. This problem becomes particularly acute when training large models on long-response tasks, leading to catastrophic model collapse. The ill-posed nature of the GRPO objective stems from a misapplication of importance sampling weights. The principle of importance sampling is to estimate the expectation of a function $f$ under a target distribution $\pi_\text{tar}$ by re-weighting samples drawn from a behavior distribution $\pi_\text{beh}$:

\begin{align}
\mathbb{E}_{z \sim \pi_\text{tar}} \left[ f(z) \right] = \mathbb{E}_{z \sim \pi_\text{beh}} \left[ \frac{ \pi_\text{tar} (z) }{ \pi_\text{beh} (z)} \, f(z) \right].
\end{align}

A key aspect of this approach is that effective correction for distribution mismatch using the importance weight $\frac{ \pi_\text{tar} (z) }{ \pi_\text{beh} (z)}$ depends on taking an average over a large number of samples ($N \gg 1$) drawn from the behavior distribution $\pi_\text{beh}$.

Conversely, GRPO introduces the importance weight $\frac{ \pi_{\theta} (y_{i,t} | x, y_{i,<t}) }{ \pi_{\theta_\text{old}} (y_{i,t} | x,y_{i,<t})}$ at each individual token position $t$. Here, the weight calculation is based on only a single token sample $y_{i,t}$ from the next-token distribution $\pi_{\theta_\text{old}}( \cdot | x, y_{i, <t})$. This methodology is inadequate for rectifying the mismatch between distributions as intended, and instead, it introduces substantial variance in the training gradients. Such variance accumulates throughout longer sequences and is worsened further by the use of clipping. Empirically, we have found that this often results in a catastrophic collapse of the model, after which recovery is typically not possible. Attempts to resume training—whether by reverting to earlier checkpoints, carefully adjusting hyperparameters (such as clipping thresholds), extending the generation horizon, or altering RL-related queries—have not reversed such collapses.

These findings reveal an inherent flaw in the design of GRPO. Specifically, the ineffectiveness of token-level importance weights underscores a central tenet: \textbf{the granularity of the optimization objective must be consistent with that of the reward signal}. Given that rewards are computed at the level of the entire sequence, applying off-policy corrections at the token level is fundamentally misaligned. This insight motivates us to abandon token-level objectives and focus on implementing importance weighting and optimization at the \textit{sequence level} instead.

\section{Algorithm}

\subsection{GSPO: Group Sequence Policy Optimization}

Although the token-level importance weight $\frac{ \pi_{\theta} (y_{i,t} | x, y_{i,<t}) }{ \pi_{\theta_\text{old}} (y_{i,t} | x,y_{i,<t})}$ poses challenges within GRPO, our findings indicate that, in the domain of language generation, the \textit{sequence-level} importance weight $\frac{ \pi_{\theta} (y | x) }{ \pi_{\theta_\text{old}} (y | x)}$ possesses a well-founded theoretical interpretation: it quantifies the extent to which a response $y$ drawn from $\pi_{\theta_\text{old}} (\cdot | x)$ diverges from the distribution under $\pi_{\theta} (\cdot | x)$. This perspective is inherently consistent with sequence-level rewards, making the importance weight a useful and interpretable gauge for the clipping strategy.

Based on this straightforward observation, we propose the \textbf{Group Sequence Policy Optimization (GSPO)} algorithm. GSPO employs the following sequence-level optimization objective:

\begin{align}
\mathcal{J}_\text{GSPO} (\theta) =
\mathbb{E}_{ x \sim \mathcal{D},\, \{y_i\}_{i=1}^G \sim \pi_{\theta_\text{old}}( \cdot | x) }
\left[ 
\frac{1}{G} \sum_{i=1}^{G}
\min \left( s_{i}(\theta)  \widehat{A}_{i},  \, \mathrm{clip} \left( s_{i}(\theta), 1 - {\varepsilon}, 1 + {\varepsilon} \right) \widehat{A}_{i} \right) 
\right],
\label{equ:gspo}
\end{align}

where we adopt the group-based advantage estimation:

\begin{align}
\widehat{A}_{i} = \frac{r(x, y_i) - \mathrm{mean} \left( \{ r(x, y_i) \}_{i=1}^G \right) }{ \mathrm{std} \left( \{ r(x, y_i) \}_{i=1}^G \right) },
\end{align}

and define the importance ratio $s_{i}(\theta)$ based on sequence likelihood \citep{click}:

\begin{align}
s_{i}(\theta) = \left( \frac{ \pi_{\theta} (y_i | x) }{ \pi_{\theta_\text{old}} (y_i | x)} \right)^{\frac{1}{|y_i|}}
=
\exp \left( \frac{1}{|y_i|} \sum_{t=1}^{|y_i|} \log \frac{ \pi_{\theta} (y_{i,t} | x, y_{i,<t}) }{ \pi_{\theta_\text{old}} (y_{i,t} | x,y_{i,<t})} \right).
\end{align}

Consequently, GSPO performs clipping at the level of complete responses rather than at the token level, thereby eliminating excessively “off-policy” samples from gradient computation. This approach aligns with both sequence-level reward mechanisms and optimization objectives. Additionally, we employ length normalization in $s_{i}(\theta)$, which serves to minimize variance and maintain $s_{i}(\theta)$ within a consistent numerical interval. Without such normalization, minor changes in the likelihood of individual tokens can lead to significant variability in the overall sequence importance ratio, and responses of differing lengths would necessitate distinct clipping intervals. It is also worth highlighting that the clipping intervals used in GSPO tend to differ by orders of magnitude from those used in prior methods (such as GRPO), largely because of the variation in how importance ratios are defined in these algorithms.

\subsection{Gradient Analysis}
\label{subsec:gradient}

We can derive the gradient of the GSPO objective as follows (clipping is omitted for brevity):

\begin{align}
\nabla_{\theta} \mathcal{J}_\text{GSPO} (\theta)
=&\ 
\nabla_{\theta} \mathbb{E}_{ x \sim \mathcal{D},\, \{y_i\}_{i=1}^G \sim \pi_{\theta_\text{old}}( \cdot | x) }
\left[ 
\frac{1}{G} \sum_{i=1}^{G} s_{i}(\theta) \widehat{A}_{i}
\right] \\
= &\ 
\mathbb{E}_{ x \sim \mathcal{D},\, \{y_i\}_{i=1}^G \sim \pi_{\theta_\text{old}}( \cdot | x) }
\left[ 
\frac{1}{G} \sum_{i=1}^{G}
s_{i}(\theta) \widehat{A}_{i}
\cdot \nabla_{\theta} \log s_{i}(\theta)
\right] \\
=&\ 
\mathbb{E}_{ x \sim \mathcal{D},\, \{y_i\}_{i=1}^G \sim \pi_{\theta_\text{old}}( \cdot | x) }
\left[ 
\frac{1}{G} \sum_{i=1}^{G} 
\left( \frac{ \pi_{\theta} (y_i | x) }{ \pi_{\theta_\text{old}} (y_i | x)} \right)^{\frac{1}{|y_i|}} \widehat{A}_{i} 
\cdot \frac{1}{|y_i|} \sum_{t=1}^{|y_i|} \nabla_{\theta} \log \pi_{\theta} (y_{i,t} | x, y_{i,<t}) 
\right].
\label{equ:gspo_grad}
\end{align}

For comparison, the gradient of the GRPO objective is as follows (note that $\widehat{A}_{i,t} = \widehat{A}_{i}$):

\begin{align}
\nabla_{\theta} \mathcal{J}_\text{GRPO} (\theta) 
=&\ 
\nabla_{\theta} \mathbb{E}_{ x \sim \mathcal{D},\, \{y_i\}_{i=1}^G \sim \pi_{\theta_\text{old}}( \cdot | x) }
\left[ \frac{1}{G} \sum_{i=1}^{G} \frac{1}{|y_i|} \sum_{t=1}^{|y_i|} 
w_{i,t}(\theta) \widehat{A}_{i,t}
\right] \\
=&\
\mathbb{E}_{ x \sim \mathcal{D},\, \{y_i\}_{i=1}^G \sim \pi_{\theta_\text{old}}( \cdot | x) }
\left[ \frac{1}{G} \sum_{i=1}^{G} \widehat{A}_{i} 
\cdot \frac{1}{|y_i|} \sum_{t=1}^{|y_i|} 
\frac{ \pi_{\theta} (y_{i,t} | x, y_{i,<t}) }{ \pi_{\theta_\text{old}} (y_{i,t} | x,y_{i,<t})} 
\nabla_{\theta} \log \pi_{\theta} (y_{i,t} | x, y_{i,<t})  
\right].
\end{align}

The principal difference between GSPO and GRPO centers on \textit{the manner in which they assign weights to the gradients corresponding to the log likelihoods of individual tokens}. GRPO employs a weighting scheme based on the so-called “importance weight,” specifically $\frac{ \pi_{\theta} (y_{i,t} | x, y_{i,<t}) }{ \pi_{\theta_\text{old}} (y_{i,t} | x,y_{i,<t})}$ for each token. These weights are not uniformly distributed; for instances with $\widehat{A}_{i} > 0$, their values range between $(0, 1+\varepsilon]$, while for $\widehat{A}_{i} < 0$, they can span $[1-\varepsilon, +\infty)$. Such variability in token weighting is significant and can intensify over the course of training, potentially introducing erratic behaviors. By comparison, GSPO applies identical weights to all tokens within a response, thus removing this source of instability inherent to GRPO.

\subsection{GSPO-token: A Token-level Objective Variant}

In scenarios like multi-turn RL, we may desire a finer-grained advantage adjustment than the sequence level. To this end, we introduce a token-level objective variant of GSPO, namely \textbf{GSPO-token}, to allow token-wise advantage customization:

\begin{align}
\mathcal{J}_\text{GSPO-token}(\theta)
=
\mathbb{E}_{ x \sim \mathcal{D},\, \{y_i\}_{i=1}^G \sim \pi_{\theta_\text{old}}( \cdot | x) }
\left[ 
\frac{1}{G} \sum_{i=1}^{G} \frac{1}{|y_i|} \sum_{t=1}^{|y_i|} 
\min \left( s_{i,t}(\theta) \widehat{A}_{i,t},  \, \mathrm{clip} \left( s_{i,t}(\theta), 1 - {\varepsilon}, 1 + {\varepsilon} \right) \widehat{A}_{i,t} \right) 
\right],
\label{equ:gspo-token}
\end{align}

where

\begin{align}
s_{i,t}(\theta) 
= \mathrm{sg} \left[ s_{i}(\theta) \right]  \cdot \frac{ \pi_{\theta} (y_{i,t} | x, y_{i,<t}) }{ \mathrm{sg} \left[ \pi_{\theta} (y_{i,t} | x, y_{i,<t}) \right] },
\end{align}

and $\mathrm{sg}[\cdot]$ denotes only taking the numerical value but stopping the gradient, corresponding to the \texttt{detach} operation in PyTorch. The gradient of GSPO-token can be derived as:

\begin{align}
\nabla_{\theta} \mathcal{J}_\text{GSPO-token}(\theta)
=&\ 
\nabla_{\theta} \mathbb{E}_{ x \sim \mathcal{D},\, \{y_i\}_{i=1}^G \sim \pi_{\theta_\text{old}}( \cdot | x) }
\left[ 
\frac{1}{G} \sum_{i=1}^{G}
\frac{1}{|y_i|} \sum_{t=1}^{|y_i|} 
s_{i,t}(\theta) \widehat{A}_{i,t}
\right] \\
=&\ 
\mathbb{E}_{ x \sim \mathcal{D},\, \{y_i\}_{i=1}^G \sim \pi_{\theta_\text{old}}( \cdot | x) }
\left[ 
\frac{1}{G} \sum_{i=1}^{G} s_i (\theta) \cdot \frac{1}{|y_i|} \sum_{t=1}^{|y_i|} 
\widehat{A}_{i,t} \frac{ \nabla_{\theta} \pi_{\theta} (y_{i,t} | x, y_{i,<t}) }{ \pi_{\theta} (y_{i,t} | x, y_{i,<t}) }
\right] \\
=&\ 
\mathbb{E}_{ x \sim \mathcal{D},\, \{y_i\}_{i=1}^G \sim \pi_{\theta_\text{old}}( \cdot | x) }
\left[ 
\frac{1}{G} \sum_{i=1}^{G}  \left( \frac{ \pi_{\theta} (y_i | x) }{ \pi_{\theta_\text{old}} (y_i | x)} \right)^{\frac{1}{|y_i|}}
\cdot \frac{1}{|y_i|} \sum_{t=1}^{|y_i|}
\widehat{A}_{i,t} \nabla_{\theta} \log \pi_{\theta} (y_{i,t} | x, y_{i,<t})  
\right].
\label{equ:gspo-token_grad}
\end{align}

It is important to observe that $\frac{ \pi_{\theta} (y_{i,t} | x, y_{i,<t}) }{ \mathrm{sg} \left[ \pi_{\theta} (y_{i,t} | x, y_{i,<t}) \right] }$ evaluates to 1, which implies that $s_{i,t}(\theta)$ shares the same numerical value as $s_{i}(\theta)$. Examining Equation~\eqref{equ:gspo} alongside~\eqref{equ:gspo-token}, as well as Equation~\eqref{equ:gspo_grad} with~\eqref{equ:gspo-token_grad}, reveals that GSPO-token and GSPO produce equivalent results in terms of the optimization target, the clipping rule, and the analytic gradient, provided that the advantage for each token within the output $y_i$ is set uniformly (that is, $\widehat{A}_{i,t} = \widehat{A}_{i}$ for all $t$). Nonetheless, GSPO-token offers greater versatility by allowing distinct advantage values to be assigned to individual tokens.

\section{Experiments and Discussion}

\subsection{Empirical Results}

We conduct experiments utilizing a cold-start configuration that has been fine-tuned from the Qwen3-30B-A3B-Base model, presenting both reward progression during training and evaluation performance curves on several benchmarks: AIME'24 (mean Pass@1 across 32 samples), LiveCodeBench (202410-202502, mean Pass@1 with 8 samples per task), and CodeForces (measured by Elo Rating). For reinforcement learning (RL) training, rollout batches are subdivided into four distinct mini-batches to apply gradient updates. Within the GSPO framework, the clipping thresholds in Equation~\eqref{equ:gspo} are assigned values of 3e-4 (left) and 4e-4 (right). Comparisons are drawn with the GRPO method, which serves as the baseline; in this case, the left and right clipping parameters in Equation~\eqref{equ:grpo} are tuned to 0.2 and 0.27, respectively, to enable equitable evaluation. It is important to observe that GRPO requires the use of the Routing Replay training approach to enable stable convergence for MoE RL; further elaboration on this is provided in \S~\ref{subsec:routing_replay}. By contrast, \textit{GSPO removes the dependency on this training technique}.

As illustrated in Figure~\ref{fig:results}, GSPO enables consistent and stable progress during training. Our findings indicate that \textbf{GSPO is capable of driving ongoing improvements in performance by increasing computational resources for training, periodically refreshing the query set, and lengthening the output sequences}. Additionally, GSPO outperforms GRPO in terms of training efficiency, attaining superior accuracy and benchmark outcomes while utilizing identical training compute and query resources. Furthermore, GSPO has been effectively integrated into RL training for the latest Qwen3 models, providing compelling evidence for GSPO's effectiveness in harnessing the scalability of RL for large language models.

\subsection{Curious Observation on Clipping Fractions}

An important difference between GSPO and GRPO lies in GSPO's method of clipping, which is applied to whole sequences rather than individual tokens. As illustrated in Figure~\ref{fig:clipping}, there exists a gap of two orders of magnitude in the proportion of clipped tokens when comparing GSPO and GRPO; notably, this difference persists regardless of changes made to the clipping ranges. Nonetheless, even though GSPO clips considerably more tokens and thus incorporates fewer tokens in its training or gradient estimation process, it still outperforms GRPO in training efficiency. This surprising outcome—where the removal of a much larger number of tokens enhances training efficiency—suggests that GRPO’s token-wise gradient estimations are intrinsically noisy and less effective at leveraging samples. Conversely, GSPO's sequence-level methodology delivers a stronger and more consistent learning signal.

\subsection{Benefit of GSPO for MoE Training}
\label{subsec:routing_replay}

\paragraph{Background}
In contrast to RL training with dense models, the inherently sparse activations of MoE models pose distinct challenges with respect to stability. Specifically, our observations reveal that the use of the GRPO algorithm can lead to considerable \textit{expert-activation volatility} within MoE architectures, which undermines the convergence of RL training. To illustrate, following one or several steps of gradient-based optimization, the set of experts selected for identical responses may shift markedly. For instance, in the case of the 48-layer Qwen3-30B-A3B-Base model, an examination after each RL gradient update shows that, for the same rollout instance, approximately 10\% of the activated experts under the revised policy $\pi_{\theta}$ differ from those selected by the previous policy $\pi_{\theta_\text{old}}$. Such variability tends to intensify in deeper MoE configurations, thereby causing sharp oscillations in the token-level importance ratios $w_{i,t}(\theta) = \frac{ \pi_{\theta} (y_{i,t} | x, y_{i,<t}) }{ \pi_{\theta_\text{old}} (y_{i,t} | x,y_{i,<t})}$ and rendering them unreliable, as elaborated in \S~\ref{sec:motivation} and~\ref{subsec:gradient}. As a consequence, the standard convergence properties of RL training are compromised.

\paragraph{Our Previous Approach} In our earlier work, we adopted the \textbf{Routing Replay} training scheme to address this issue. This method involves storing the set of experts activated by $\pi_{\theta_\text{old}}$ and subsequently enforcing these cached routing patterns during evaluation of $\pi_{\theta}$ in the computation of the importance ratios $w_{i,t}(\theta) = \frac{ \pi_{\theta} (y_{i,t} | x, y_{i,<t}) }{ \pi_{\theta_\text{old}} (y_{i,t} | x,y_{i,<t})}$. By aligning the network activation for both $\pi_{\theta} (y_{i,t} | x, y_{i,<t})$ and $\pi_{\theta_\text{old}} (y_{i,t} | x,y_{i,<t})$ for each token $y_{i,t}$, this strategy stabilizes token-level importance ratios and guarantees that the same subset of experts is optimized throughout the training steps. As shown in Figure~\ref{fig:routing_replay}, Routing Replay proves vital for ensuring the successful convergence of MoE model GRPO training.

\paragraph{Benefit of GSPO} While Routing Replay allows MoE models to achieve proper GRPO training convergence, its method of reusing routing patterns increases both memory usage and communication costs, and further constrains the MoE's potential capacity. Alternatively, GSPO, as depicted in Figure~\ref{fig:results}, operates independently from Routing Replay. It reliably calculates the importance ratios $s_i(\theta)$ in the standard manner, converges as expected, and optimizes steadily. The central observation is that GSPO emphasizes sequence likelihood (i.e., $\pi_{\theta} (y_i | x)$) rather than being affected by the likelihood of individual tokens (i.e., $\pi_{\theta} (y_{i,t} | x, y_{i,<t})$). Since the MoE architecture consistently upholds its core language modeling abilities, the sequence likelihood remains largely stable. Ultimately, GSPO addresses the instability in expert activation inherent to MoE models, thereby removing the necessity for sophisticated techniques such as Routing Replay. As a result, the approach not only streamlines and reinforces the training regimen but also ensures that the model accesses its entire capacity without imposed limitations.

\subsection{Benefit of GSPO for RL Infrastructure}

Due to the differences in numerical precision between engines utilized for training, such as Megatron, and those designed for inference, like SGLang and vLLM, it is common in practice to recalculate the likelihoods of responses sampled from the previous policy $\pi_{\theta_\text{old}}$ using the training engine. In contrast, GSPO relies on optimizing based on sequence-level, not token-level, likelihoods; sequence-level likelihoods tend to be significantly less sensitive to precision variation. As a result, GSPO enables direct utilization of likelihood outputs from the inference engine during optimization, obviating the need for recomputation with the training engine. This capability offers particular advantages in applications including partial rollout, multi-turn RL, and systems that decouple training from inference.

\section{Conclusion}

We introduce Group Sequence Policy Optimization (GSPO), a novel algorithm in reinforcement learning tailored for the training of large language models. Leveraging the core concepts of importance sampling, GSPO establishes sequence-level importance ratios grounded in sequence likelihood, and applies clipping, reward assignment, and optimization at the sequence level. Empirical results reveal that GSPO surpasses GRPO in training stability, efficiency, and overall performance, demonstrating distinct advantages in large-scale RL training for MoE models and underpinning substantial enhancements seen in the recent Qwen3 models. As GSPO provides a robust and scalable methodological foundation, our ongoing efforts aim to further scale reinforcement learning, anticipating significant progress in intelligence as a consequence.

\end{document}